\begin{abstract}

\noindent\textbf{Purpose:}
The apparent diffusion coefficient (ADC) is the de facto biomarker of prostate diffusion MRI, yet the multi-compartment models invoked to explain it are not measured directly. Bayesian spectral decomposition of extended b-value prostate DWI was used to identify which compartments are recoverable and whether they explain ADC's near-optimality for region-of-interest (ROI) tumor detection.

\noindent\textbf{Theory and Methods:}
Multi-b DWI (15 b-values to 3500~s/mm$^2$) from 56 patients (149 ROIs: 40 tumor, 109 normal) was decomposed into eight diffusivity components by constrained-ridge maximum a posteriori (MAP) and fully Bayesian No-U-Turn Sampler (NUTS) inference, the latter also inferring per-ROI noise and component uncertainty. Tumor detection used zone-specific leave-one-out logistic regression with patient-level bootstrap confidence intervals, comparing ADC against spectral models from single fractions to all eight.

\noindent\textbf{Results:}
Only restricted diffusion and free water were well identified (posterior coefficient of variation $0.20$ and $0.32$); intermediate components were small and weakly constrained, even under wider priors. The eight-fraction classifier matched ADC (AUC $0.93$/$0.92$ vs.\ $0.95$/$0.98$, peripheral/transition zones; overlapping bootstrap CIs); the restricted and free-water fractions alone matched it ($0.93$/$0.94$), whereas the five intermediate fractions alone fell to $0.82$/$0.80$. ADC was near-perfectly anti-correlated with the optimal spectral discriminant ($|r| > 0.97$). MAP and NUTS agreed, the posterior adding no diagnostic information.

\noindent\textbf{Conclusion:}
For ROI-level detection the spectrum is diagnostically equivalent to ADC, which it mechanistically explains rather than supersedes---detection rides one restricted-to-free-water axis. The Bayesian treatment quantifies what the acquisition can resolve: the two extreme compartments well, the intermediate ones weakly; grade-dependent shifts mark where the spectrum may add value.

\end{abstract}
